\documentclass[12pt,a4paper]{article}
\usepackage[a4paper,left=2.5cm,right=2cm,top=2cm,bottom=2cm]{geometry}
\usepackage{fontspec}
\setmainfont{DejaVu Serif}
\setmonofont{DejaVu Sans Mono}
\usepackage{setspace}
\onehalfspacing
\usepackage{enumitem}
\usepackage{booktabs}
\usepackage{tabularx}
\usepackage{array}
\usepackage{xcolor}
\usepackage{listings}
\usepackage{hyperref}
\hypersetup{hidelinks}
\usepackage{titlesec}
\titleformat{\section}{\normalfont\Large\bfseries}{\thesection.}{0.5em}{}
\titleformat{\subsection}{\normalfont\large\bfseries}{\thesubsection.}{0.5em}{}
\setlength{\parindent}{1.25cm}
\setlength{\parskip}{0pt}
\renewcommand{\contentsname}{Contents}
\definecolor{codebg}{RGB}{246,246,246}
\definecolor{codeframe}{RGB}{215,215,215}
\lstset{basicstyle=\ttfamily\small,backgroundcolor=\color{codebg},frame=single,
  rulecolor=\color{codeframe},breaklines=true,columns=fullflexible,keepspaces=true,
  showstringspaces=false,aboveskip=0.8em,belowskip=0.5em}
\newcommand{\figcap}[1]{\begin{center}\small\textit{#1}\end{center}\vspace{0.4em}}

\begin{document}
\begin{titlepage}
\begin{center}
{\large Ministerul Educației și Cercetării al Republicii Moldova}\\[0.3cm]
{\large Universitatea Tehnică a Moldovei}\\[0.3cm]
{\large Facultatea Calculatoare, Informatică și Microelectronică}\\[0.3cm]
{\large Departamentul Ingineria Software și Automatică}\\[0.3cm]
{\large Program: Software Engineering}\\[2.0cm]
{\Large \textbf{Operating Systems}}\\[0.8cm]
{\LARGE \textbf{Laboratory Work No. 7}}\\[0.5cm]
{\Large \textbf{Three Small Programs: Processes, Redirection, a Pipe}}\\[0.35cm]
{\large Track A: The Real OS, from the Outside (Linux)}
\vfill
\begin{flushright}
\begin{tabular}{ll}
\textbf{Done by:} & st. gr. FAF-241, Ermicev Sofia \\
\\
\textbf{Checked by:} & Patricia Reițman
\end{tabular}
\end{flushright}
\vfill
{\large Chișinău -- 2026}
\end{center}
\end{titlepage}
\tableofcontents
\newpage

\section{Purpose of the Work}
The purpose of this laboratory work is to reproduce basic shell mechanisms directly in C: creating a process with \texttt{fork()}, replacing it with \texttt{exec()}, waiting for it, observing a zombie, redirecting standard output with file descriptors, and connecting two processes through a pipe.

\section{Part 1. fork + exec + wait}
\begin{lstlisting}
$ gcc -Wall -Wextra -o run run.c && ./run
parent: my PID is 41031
child: my PID is 41032, my parent is 41031
total 24
-rwxr-xr-x 1 sofia sofia 16400 Oct 19 14:10 run
-rw-r--r-- 1 sofia sofia   919 Oct 19 14:09 run.c
parent: child 41032 ended with status 0
\end{lstlisting}
\figcap{Figure 1.1 -- Actual execution of run with ls -l}

\begin{lstlisting}
parent: my PID is 41555
child: my PID is 41556, my parent is 41555
Mon Oct 19 14:11:58 EEST 2026
parent: child 41556 ended with status 0
\end{lstlisting}
\figcap{Figure 1.2 -- Actual execution after changing exec to date}

\begin{lstlisting}
parent: my PID is 42078
child: my PID is 42079, my parent is 42078
exec: No such file or directory
parent: child 42079 ended with status 127
\end{lstlisting}
\figcap{Figure 1.3 -- Actual failed exec experiment}

\subsection{Observe answers}
\begin{enumerate}[label=\arabic*.]
\item \texttt{perror("exec")} runs only after an unsuccessful \texttt{execlp()}, because successful \texttt{exec} replaces the old process image and does not return.
\item The child reported the PID of its actual parent; for example PID 41032 reported parent PID 41031.
\item Successful \texttt{ls} and \texttt{date} runs returned status 0. The failure path explicitly used \texttt{exit(127)}, so the parent observed 127.
\end{enumerate}

\section{Part 2. A Zombie, on Purpose}
\begin{lstlisting}
$ ./zombie > /dev/null &
$ sleep 2
$ ps -o pid,ppid,stat,cmd
    PID    PPID STAT CMD
  40200   40193 Ss   -bash
  44268   40200 S    ./zombie
  44270   44268 Z    [ls] <defunct>
  44271   40200 R+   ps -o pid,ppid,stat,cmd
\end{lstlisting}
\figcap{Figure 2.1 -- Actual zombie process shown by ps}

\begin{lstlisting}
$ ps -o pid,ppid,stat,cmd
    PID    PPID STAT CMD
  40200   40193 Ss   -bash
  44456   40200 R+   ps -o pid,ppid,stat,cmd
[1]+  Done                    ./zombie > /dev/null
\end{lstlisting}
\figcap{Figure 2.2 -- Actual process table after waitpid collected the child}

\subsection{Observe answers}
\begin{enumerate}[label=\arabic*.]
\item \texttt{STAT=Z} means the child has terminated but its parent has not yet collected its status.
\item After the parent's 20-second delay, \texttt{waitpid()} collected the child and the \texttt{<defunct>} entry disappeared.
\end{enumerate}

\section{Part 3. Redirection: ls -l > out.txt}
\begin{lstlisting}
$ gcc -Wall -Wextra -o redirect redirect.c && ./redirect
open gave descriptor 1
done: the listing is in out.txt
\end{lstlisting}
\figcap{Figure 3.1 -- Actual redirect program output}

\begin{lstlisting}
$ head -3 out.txt
total 68
-rw-r--r-- 1 sofia sofia     0 Oct 19 14:26 out.txt
-rwxr-xr-x 1 sofia sofia 16344 Oct 19 14:26 redirect
\end{lstlisting}
\figcap{Figure 3.2 -- Actual beginning of the redirected ls output}

\subsection{Observe answers}
\begin{enumerate}[label=\arabic*.]
\item \texttt{open()} returned descriptor 1 because \texttt{close(1)} made it the smallest available descriptor.
\item \texttt{ls} still wrote to stdout; only the object referenced by descriptor 1 changed from the terminal to \texttt{out.txt}.
\item The same method can redirect stderr by closing descriptor 2 before opening the error file.
\end{enumerate}

\section{Part 4. A Pipe: ls /etc | wc -l}
\begin{lstlisting}
$ gcc -Wall -Wextra -o pipeline pipeline.c && ./pipeline
183
$ ls /etc | wc -l
183
\end{lstlisting}
\figcap{Figure 4.1 -- Actual manual and shell pipeline results}

\begin{lstlisting}
$ gcc -Wall -Wextra -o pipeline pipeline.c
$ ./pipeline
^C
\end{lstlisting}
\figcap{Figure 4.2 -- Actual blocked pipeline when the parent kept the write end open}

\subsection{Observe answers}
\begin{enumerate}[label=\arabic*.]
\item The manually constructed pipeline and the shell pipeline both produced 183.
\item When the parent deliberately kept its copy of the write end open, the \texttt{wc} child did not terminate.
\item \texttt{wc} was waiting for EOF; EOF cannot be delivered while any process still has a write descriptor referring to that pipe.
\end{enumerate}

\section{Conclusion}
The programs reproduced core shell behavior using Linux system calls. The work demonstrated process creation and replacement, child-status collection, the zombie state, descriptor-based redirection, and the need to close unused pipe ends so readers can eventually receive EOF.

\end{document}
